\documentclass[11pt,a4paper]{article}
\usepackage[margin=25mm]{geometry}
\usepackage{amsmath,amssymb,amsthm,hyperref}
\setlength{\emergencystretch}{2em}
\newtheorem{lemma}{Lemma}
\newtheorem{proposition}[lemma]{Proposition}
\newcommand{\gen}[1]{\langle#1\rangle}
\title{An exact residual case under an alternating tail}
\author{Human Compression project: bounded L044 proof component}
\date{}
\begin{document}\maketitle
Retain the hypotheses and equality parameters of the numerical
intransitive-side support theorem. Thus $G\leq A_n$ is a transitive
connected string C-group, all maximal omissions are intransitive,
and a unique eligible interior $i$-edge $ab$ joins omission
orbits $A,B$. Here $i=k-1\geq1$, $h=i+1=k$, and
$A=U\sqcup V$ consists of two faithful natural $S_k$-orbits
of $L=G_{<i}$, matched by $\tau=\rho_{i+1}$.
The group $L$ fixes $B$ pointwise.
The tail $K=G_{>h}$ fixes exactly $A\cup\{b\}$.
Assume additionally that $K$ is the full natural alternating
group on $X=B\setminus\{b\}$.
Write $U=\{u_1,\ldots,u_k\}$, $V=\{v_1,\ldots,v_k\}$,
with $a=u_k$ and $\tau$ exchanging $u_j,v_j$ for every $j$.
Put $\lambda=\prod_j(u_j\ v_j)$ and $s=\rho_i$.

\begin{lemma}[The crossing under the diagonal omission]\label{crossing}
For $k\geq4$, evenness forces $k$ even and
\[
                   s=(u_k\ b)\prod_{j<k}(u_j\ v_j).       \tag{1}
\]
The case $k=3$ is impossible. For $k=2$, formula (1) also holds.
\end{lemma}
\begin{proof}
The tail is nontrivial and generated by involutions, so
$q=|X|\geq5$: the involutions of $A_4$ generate only its
Klein four subgroup, and the smaller alternating groups do
not supply a nontrivial involutory generating family.
The generator $s$ commutes with $K$, hence preserves its sole
nontrivial orbit $X$ and fixes it pointwise. The natural
alternating action has trivial centralizer. Thus
$s=(u_k\ b)w$, with $w$ supported on $A\setminus\{u_k\}$.

Let $H=G_{<i-1}$. Its $k-2$ generators fix $u_k$ and $v_k$.
Their faithful images in $S_{k-1}$ are independent, so Whiston's
equality clause identifies $H$ as diagonal $S_{k-1}$ on the
remaining two rows. For $k\geq3$ these are the only two fixed
points of $H$ in $A$. String commutation makes $w$ centralize
$H$; since it fixes $u_k$, it also fixes $v_k$.

For $k\geq4$, the centralizer on the remaining rows is $C_2$,
generated by their full exchange. Indeed a commuting permutation
permutes the two natural $S_{k-1}$-orbits, and the natural action
has trivial centralizer. Evenness of $s$ forces $w$ to be the
exchange, and requires its $k-1$ transpositions to be odd.
This proves (1) and the evenness of $k$.

For $k=3$, label the two distinguished $L$ generators
$\rho_0=(u_1\ u_2)(v_1\ v_2)$ and
$\rho_1=(u_2\ u_3)(v_2\ v_3)$, as their faithful $S_3$ action
and the fixed anchor permit. An odd involution $w$ on the other
four points centralizing $\rho_0$ must be one of
$(u_1\ u_2)$ and $(v_1\ v_2)$. The group $G_0$ contains
$G_{>2}$, which is transitive on $B$ and exchanges every pair
$\{u_j,v_j\}$, as well as $\rho_1$ and $s$.
The permutation $\rho_1$ joins pairs 2 and 3, $w$ joins pair 1
to pair 2, and the cross-edge of $s$ joins them to $B$.
Hence $G_0$ is transitive, a contradiction.
Finally, for $k=2$, the only possible odd involution on the
three points $A\setminus\{u_2\}$ is a transposition, but
$w$ need not initially fix $v_2$. If it exchanges $v_2$ with
one of $u_1,v_1$, then $G_0=\gen{s,G_{>1}}$ is transitive:
the latter factor exchanges each vertical pair and is transitive
on $B$, while $s$ joins the two pairs and $B$. Thus intransitivity
forces $w=(u_1\ v_1)$, proving (1) in this case too.
\end{proof}

\begin{lemma}[A vertical vector from a five-cycle]\label{vector}
Assume (1). If $g=(u_k\ u_j)(v_k\ v_j)\in L$ for $j\ne k$,
then
\[
                     (s g)^5=\prod_{l\notin\{j,k\}}(u_l\ v_l).
                                                               \tag{2}
\]
For every even $k\geq4$, the group $P=\gen{L,s}$ therefore
contains $\lambda$.
\end{lemma}
\begin{proof}
On the five points $u_k,v_k,u_j,v_j,b$, the product $s g$ is
a five-cycle. On each other vertical pair it is a transposition.
These supports are disjoint, proving (2). Conjugating by diagonal
$S_k$ gives the product of all vertical flips except any specified
two. Multiplying the products which omit $\{a,c\}$ and
$\{b,c\}$ gives the flip of just pairs $a,b$, for distinct
$a,b,c$. All two-pair flips therefore belong to $P$.
For even $k$ their products include the flip of all pairs,
namely $\lambda$.
\end{proof}

\begin{proposition}[Reduction to a dihedral prefix]
Under the stated equality and alternating-tail hypotheses,
necessarily $k=2$, $i=1$, $h=2$, and $G_{<2}$ is exactly the
degree-five dihedral marking
\[
 \rho_0=(u_1\ u_2)(v_1\ v_2),\qquad
 \rho_1=(u_2\ b)(u_1\ v_1),\qquad |G_{<2}|=10.
\]
This conclusion does not assert that this residual case is
realizable with the full string intersection property.
\end{proposition}
\begin{proof}
By Lemma~\ref{crossing}, only even $k$ remain. On $A$, the
group $R=G_{>i}=\gen{\tau,K}$ acts through $\gen\lambda$.
Since $k$ is even, $\tau|_B$ is even; it moves $b$ into $X$.
The groups $K=A_X$ and $K^\tau=A_{B\setminus\{b^\tau\}}$
generate $A_B$: their natural alternating supports overlap
in at least four points, and the elementary three-cycle gluing
argument gives the alternating group on their union.
Thus $R$ is a subgroup of $\gen\lambda\times A_B$ projecting
onto both factors. If it did not contain $(\lambda,1)$, it
would be the graph of a nontrivial homomorphism $A_B\to C_2$.
Perfection of $A_B$ excludes that possibility, so $\lambda\in R$.

For $k\geq4$, Lemma~\ref{vector} also gives $\lambda\in P$.
But $P=G_{\leq i}$ and $R=G_{>i}$ have disjoint index sets;
full IP requires $P\cap R=1$. This contradiction excludes every
such $k$. At $k=2$, the two displayed involutions have product
a five-cycle, so generate precisely $D_{10}$ on the five
displayed points. They fix $X$ pointwise, proving the stated
exact prefix and nothing about its possible completion.
\end{proof}

\paragraph{Source boundary and attribution.}
This supplies an explicit group-theoretic reduction for the
intransitive-side equality step used in L044 Proposition~5.26,
physical page~34: Cameron--Fernandes--Leemans--Mixer,
\emph{Highest rank of a polytope for $A_n$}, arXiv:1605.09173v1.
It reuses the accepted numerical equality,
unique-crossing and Whiston inputs. The vertical-vector word
and the disjoint-parabolic IP obstruction are proved here.
The surviving $D_{10}$ branch remains a principal obligation;
it is not silently discarded by the source's assertion about
vertical edges. This note neither proves nor disproves the
complete Proposition~5.26 or the main alternating-group rank
theorem. The transitive-side equality and all later source
branches remain separate. No finite computation, new formal
declaration, whole-original or question credit is asserted.
\end{document}
